% !PRINTED 776
\weberpage{776}
%<<CONT
in common with $\Phi(x)$, $\Phi(x)$ must be divisible by $f_n(\alpha^{-1}x)$, and we can accordingly set, for an indeterminate $x$,

\weq{18}{\Phi(x) = f_n(\alpha^{-1}x)\ \Psi(x,\alpha)}

where $\Psi(x,\alpha)$ is an integral rational function of $x$.

If $\alpha, \alpha_1, \alpha_2, \ldots$ are all the primitive $m$th roots of unity, which, as we have assumed, are roots of an \gesperrt{irreducible} rational equation

\[
F_m(x) = 0
\]

then, by this assumption, $\alpha$ can be replaced in the identity (18) by any other one of $\alpha_1, \alpha_2, \ldots$, and it follows that $\Phi(x)$ is divisible by each of the factors

\[
f_n(\alpha^{-1}x),\ f_n(\alpha_1^{-1}x),\ f_n(\alpha_2^{-1}x) \ldots
\]

But two of these factors can have no common divisor; for if, say,

\[
f_n(\alpha^{-1}x) = 0,\ f_n(\alpha_1^{-1}x) = 0
\]

had a common root $x$, then two primitive $n$th roots of unity $r, r_1$ would have to exist such that

\[
x = \alpha r = \alpha_1 r_1
\]

held; so we should have $\alpha^n = \alpha_1^n$, and consequently, since $n$ is relatively prime to $m$, $\alpha = \alpha_1$.

Accordingly $\Phi(x)$ is divisible by the product

\[
F_{mn}(x) = f_n(\alpha^{-1}x)\ f_n(\alpha_1^{-1}x)\ f_n(\alpha_2^{-1}x) \ldots,
\]

whose roots are all the $\beta$, and this function is therefore also indecomposable, as was to be proved.

\begin{center}
\rule{4cm}{0.4pt}
\end{center}
